\section{Primeri sistema}

U ovom poglavlju biće prikazani konkretni primeri distribuiranih sistema koji koriste DAG strukturu za organizaciju transakcija i postizanje konsenzusa. Fokus će biti na projektima koji su implementirali različite pristupe DAG konsenzusu. Analiziraće se sitemi IOTA, \textit{Hedera}, \textit{Nano} i \textit{Constellation Network}.

\subsection{IOTA}

IOTA je jedna od prvih kriptovaluta koja je u potpunosti zasnovana na DAG strukturi koja, za razliku od tradicionalnih blokčejn mreža nema rudare.

U originalnoj verziji IOTA je bila zasnovana na DAG strukturi, poznatoj kao \textit{Tangle}, u kojoj je svaka nova transakcija morala da izvrši malu \textit{Proof of Work} operaciju i potvrdi dve prethodne transakcije. Izbor koje transakcije će biti potvrđene se vršio pomoću Monte Karlo metode zasnovane na Markovljevim lancima (engl. \textit{Markov Chain Monte Carlo}), koji probabilistički bira grane grafa sa najvećom akumuliranom težinom. Centralni čvor, takozvani \enquote{koordinator}, je periodično izvršavao transakcije kako bi doneo konačnu odluku koje transakcije su sigurno validne~\cite{blockeden DAG}. 

Prethodni pristup je bio izložen kritikama zbog centralizacije i u novijoj IOTA 2.0 verziji zamenjen je potpunim decentralizovanim konsenzusom koji kombinuje glasanje na samoj mreži sa grupom validator čvorova. Validatori se biraju putem dPoS \textit{(Delegated Proof-of-Stake)} mehanizma i učestvuju u donošenju odluka primenom vizantijsko-otpornih (BFT) protokola, čime se postiže deterministička finalnost transakcija. Kreiranu transakciju klijent prosleđuje \enquote{punom} čvoru, koji obavlja osnovne provere ispravnosti i distribuira je validatorima. Validatori nezavisno proveravaju transakciju i, ukoliko je validna, potpisuju je i šalju potvrdu klijentu. Nakon potpisivanja transakcije od strane dve trećine validatora, formira se sertifikat transakcije i šalje validatorima. Sertifikat se uključuje u DAG-bazirani mehanizam sekvenciranja, u kome validatori proveravaju njegovu ispravnost, izvršavaju transakciju i uspostavljaju globalni redosled transakcija. Klijenti prikupljaju potvrde o izvršenoj transakciji od validatora i nakon dve trećine prikupljenih potvrda imaju dokaz o finalizaciji transakcije~\cite{iota docs}.

Ovakav pristup omogućava zaštitu od dvostruke potrošnje, kao i horizontalnu skalabilnost. Kapacitet mreže se povećava sa rastom broja validatora i paralelne obrade transakcija, što IOTA čini posebno pogodnom za primene u internetu stvari (engl. \textit{Internet of Things}), kao što su mikroplaćanja između uređaja i praćenje senzorskih podataka. Pored toga, IOTA podržava različite modele izvršavanja pametnih ugovora uz kompatibilnost sa \textit{Ethereum} virtuelnom mašinom. U pogledu troškova, IOTA uvodi mehanizam sponzorisanih transakcija, gde treća strana (sponzor) pokriva troškove u ime korisnika, čime se olakšava korišćenje mreže i poboljšava pristupačnost aplikacija~\cite{blockeden DAG}.

% Prethodni pristup je bio izložen kritikama zbog centralizacije i u novijoj IOTA 2.0 verziji zamenjen je potpunim decentralizovanim konsenzusom koji kombinuje glasanje na samoj mreži sa grupom validatora. Svaka nova transakcija automatski daje glas transakcijama koje referencira. Posebnu ulogu imaju validator čvorovi, odabrani putem dPoS \textit{(Delegated Proof-of-Stake)} protokola, koji emituju validacione blokove sa težinski ponderisanim glasovima. Transakcija se smatra potvrđenom kada sakupi dovoljno ukupne težine glasova, tipično više od dve trećine ukupnog uloga validatora~\cite{blockeden DAG}.


% IOTA je jedna od prvih kriptovaluta koja je u potpunosti zasnovana na DAG stukturi, poznatoj kao \textit{Tangle}. Za razliku od tradicionalnih blokčejn mreža, nema rudare, već svaka nova transakcija mora da izvrši malu \textit{Proof of Work} operaciju i potvrdi dve prethodne transakcije. Izbor koje transakcije će biti potvrđene se često vrši pomoću Monte Karlo metode zasnovane na Markovljevim lancima (engl. \textit{Markov Chain Monte Carlo}), koji probabilistički bira grane grafa sa najvećom akumuliranom težinom~\cite{blockeden DAG}. 

% U originalnoj IOTA mreži, centralni čvor, takozvani \enquote{koordinator}, je periodično izvršavao transakcije kako bi doneo konačnu odluku koje transakcije su sigurno validne. Ovaj pristup je bio izložen kritikama zbog centralizacije i u novijoj IOTA 2.0 verziji zamenjen je potpunim decentralizovanim konsenzusom koji kombinuje glasanje na samoj mreži sa grupom validatora. Svaka nova transakcija automatski daje glas transakcijama koje referencira. Posebnu ulogu imaju validator čvorovi, odabrani putem dPoS \textit{(Delegated Proof-of-Stake)} protokola, koji emituju validacione blokove sa težinski ponderisanim glasovima. Transakcija se smatra potvrđenom kada sakupi dovoljno ukupne težine glasova, tipično više od dve trećine ukupnog uloga validatora~\cite{blockeden DAG}.

%\colorbox{yellow}{IOTA koristi dPoS i Mysticeti protokole}

\subsection{Hedera}

\textit{Hedera} je javna distribuirana glavna knjiga koja, za razliku od većine javnih blokčejn mreža, ima dozvoljeno upravljanje (engl. \textit{permissioned governance}), što znači da samo odobreni članovi Hedera saveta mogu da upravljaju čvorovima koji učestvuju u konsenzusu.

\textit{Hedera} koristi \textit{HashGraph} distribuirani konsenzus algoritam koji omogućava učesnicima mreže da se saglase o redosledu i vremenskim trenucima transakcija bez oslanjanja na rudarenje ili centralnog lidera. \textit{HashGraph} koristi DAG u kome svaki čvor razmenjuje informacije sa drugim nasumično odabranim čvorovima putem \textit{\enquote{Gossip About Gossip}} protokola. To znači da komunikacijom prosleđuje ne samo transakcije, već i podatke o istoriji same komunikacije između čvorova, to jest, događaje. Na osnovu ovakvog grafa, svaki čvor može lokalno da rekonstruiše kompletnu istoriju širenja informacija. Radi postizanja konsenzusa o redosledu događaja, čvor koristi virtuelno glasanje. Bez slanja posebnih glasova kroz mrežu, izborom odabranih događaja koji su brzo i široko distribuirani (\textit{famous witnesses}), čvor lokalno određuje redosled događaja. Ovakav pristup postiže asinhronu vizantijsku toleranciju na greške (aBFT - \textit{Asynchronous Byzantine Fault Tolerance}), što znači da sistem može da toleriše do jedne trećine zlonamernih čvorova, dok asinhronost podrazumeva nezavisnost od fiksnih vremenskih prozora i garantuje postizanje konsenzusa čak i u uslovima nepredvidljivih mrežnih kašnjenja~\cite{baird HashGraph,coin HashGraph}.

% Ovakav pristup postiže asinhronu vizantijsku toleranciju na greške (aBFT - \textit{Asynchronous Byzantine Fault Tolerance}), što znači da sistem može da izdrži zlonamerne čvorove sve dok ih je manje od jedne trećine~\cite{baird HashGraph,coin HashGraph}. 

Sistem zbog ovih karakteristika omogućava visok protok transakcija, nisku latenciju i pravedno potvrđivanje transakcija bez potrebe za energetski intenzivnim procesima~\cite{baird HashGraph,coin HashGraph}. Mreža omogućava obradu do 10.000 transakcija u sekundi (TPS) uz niske i predvidljive transakcione naknade. Pogodna je za primenu u različitim industrijama i ima podršku za pametne ugovore kompatibilne sa \textit{Ethereum} virtuelnom mašinom~\cite{hedera}.

\subsection{Nano}

\textit{Nano} je kriptovaluta koja koristi \textit{block-lattice} strukturu za organizaciju transakcija. U \textit{Nano} mreži svaki čvor ima svoj lanac transakcija, a transferi između čvorova se sprovode pomoću slanja i prijema blokova na odgovarajućim lancima.

Konsenzus u \textit{Nano} mreži se ostvaruje putem \textit{Open Representative Voting} (ORV) mehanizma. Svaki korisnik mreže može delegirati svoje glasove (engl. \textit{voting weight}), proporcionalne stanju na njegovom računu, predstavničkim čvorovima (engl. \textit{representatives}). Predstavnički čvorovi su čvorovi koji su konstantno pristutni na mreži i glasaju o transakcijama koje prime. Predstavnički čvorovi koji sakupe dovoljno delegiranog biračkog prava, postaju glavni predstavnici (engl. \textit{Principal Representative}) i njihovi glasovi se dodatno prenose kroz mrežu radi bržeg postizanja konsenzusa. Konsenzus se postiže kada neki blok dobije podršku više od dve trećine ukupnog glasačkog prava na mreži, nakon čega svaki čvor nezavisno potvrđuje i \enquote{cementira} trasakciju kao nepovratnu. Za razliku od klasičnih \textit{Proof-of-Stake} sistema, ORV ne zahteva zaključavanje sredstava niti generisanje blokova od strane validatora, a korisnici zadržavaju kontrolu nad svojim sredstvima i mogu promeniti predstavnika u bilo kom trenutku~\cite{ORV consensus}. 

Ovaj pristup omogućava brzu, energetski efikasnu i decentralizovanu potvrdu transakcija bez naknada, jer ne postoji ekonomska nagrada za učesnike u konsenzusu. Prosečno vreme finalizacije transakcija u uobičajenim okolnostima mreže iznosi 0,2 sekunde. Mreža je fokusirana na digitalni novac i pogodna je za mikroplaćanja, kao i svakodnevne digitalne transakcije.~\cite{ORV consensus}

\subsection{Constellation Network}

\textit{Constellation Network} je distribuirani sistem zasnovan na DAG strukturi, koji koristi paralelnu validaciju podataka kroz više slojeva konsenzusa, čime se izbegavaju uska grla karakteristična za linearne lance. Osnovu sistema čini arhitektura poznata kao \textit{Hypergraph Transfer Protocol} (HGTP), koja omogućava nezavisnu obradu podataka i njihovu kasniju globalnu finalizaciju~\cite{constellation docs}.

Arhitektura \textit{Constellation Network}-a je hijerarhijski organizovana i sastoji se od dva sloja: \textit{Layer 1} (L1) i \textit{Layer 0} (L0). Na sloju L1 vrši se inicijalna validacija transakcija korišćenjem DAG strukture. Ovaj sloj može biti specijalizovan za različite tipove aplikacija ili podataka, uključujući takozvane \textit{metagraph}-ove, koji predstavljaju aplikacione podmreže sa sopstvenim pravilima validacije i logikom izvršavanja~\cite{constellation whitepaper}.

Nakon inicijalne validacije na L1 sloju, podaci se agregiraju u strukture poznate kao \textit{snapshots} i prosleđuju ka sloju L0, koji predstavlja centralni konsenzusni sloj mreže. L0 je zadužen za finalnu verifikaciju i objedinjavanje stanja svih L1 podataka i \textit{metagraph}-ova u jedinstvenu globalnu istoriju~\cite{constellation whitepaper}.

Konsenzus u \textit{Constellation Network}-u se ostvaruje putem mehanizma \textit{Proof of Reputable Observation} (PRO). Za razliku od tradicionalnih blokčejn konsenzus mehanizama, PRO se ne oslanja isključivo na računarsku snagu ili količinu uloženih sredstava, već na reputaciju validator čvorova. Validator čvorovi učestvuju u posmatranju i validaciji mrežnih događaja, pri čemu se njihova reputacija formira na osnovu istorijske tačnosti njihovog ponašanja u mreži. Reputacija direktno utiče na težinu njihovih potvrda tokom procesa konsenzusa, pri čemu sistem favorizuje pouzdane validatore, dok istovremeno umanjuje uticaj zlonamernih učesnika~\cite{constellation whitepaper,crypto PRO constellation}.

Validator čvorovi učestvuju u konsenzusu na oba sloja i nagrađuju se u tokenu mreže, nazvanom \$DAG, koji ima centralnu ulogu u ekonomiji sistema i koristi se za podsticaj validatora~\cite{constellation docs}.

Zahvaljujući navedenim karakteristikama, \textit{Constellation Network} je posebno pogodan za primene koje zahtevaju visoku propusnost, nisku latenciju i obradu velikih količina podataka, kao što su sistemi za analizu podataka, IoT infrastrukture i kompleksne Web3 aplikacije~\cite{constellation whitepaper}.